\clearpage


\addcontentsline{toc}{chapter}{Abstract (Hebrew)}
\markboth{Abstract (Hebrew)}{Abstract (Hebrew)}

\begin{otherlanguage}{hebrew}
\begin{justify}
\begin{flushright}
\label{sec:abstract-hebrew}
\chapter*{\texthebrew{תקציר}}
% \textbf{ארכיטקטורות למידה עמוקה לתקשורת ממסר דו-קפיצתית: מחקר השוואתי של אסטרטגיות ממסר קלאסיות ומבוססות רשתות נוירונים}

% חיבור זה מציג מחקר השוואתי של אסטרטגיות ממסר (relay) קלאסיות ומבוססות רשתות נוירונים עבור תקשורת שיתופית דו-קפיצתית. המחקר כולו מתבצע על תצורה קנונית אחת, הקבועה מראש ואינה משתנה לאורך קו הטיעון המרכזי: קישור SISO בשני הקפיצות, מעל דעיכת Rayleigh מהירה בלתי-תלויה ושווה-התפלגות, בבסיס-רוחב מרוכב, אפנון BPSK, ושיעור שגיאות סיביות (BER) בלתי-מקודד, כאשר פונקציית הממסר היא הציר היחיד המשתנה. נבדקות תשע שיטות ממסר: שתי גישות קלאסיות : הגברה-והעברה (AF) ופענוח-והעברה (DF) : ושבע שיטות מבוססות למידה עמוקה: למידה מפוקחת (MLP, ממסר Hybrid המתאים את עצמו לרמת ה-SNR), מודלים גנרטיביים (VAE, CGAN), וארכיטקטורות רצפים (Transformer, Mamba S6, Mamba-2 SSD). כל התוצאות מבוססות על סימולציית מונטה קרלו (100,000 ביטים לכל נקודת SNR, עשרה ניסויים, רווחי סמך של 95\%).

% בתצורה הקנונית, ממסרי הרשת הטובים ביותר עולים על AF ב-SNR נמוך ($0$–$4$ dB) אך רק משתווים ל-DF הסימבולי, אשר שולט בתחום SNR בינוני-גבוה עם אפס פרמטרים. מחקר מורכבות חושף יחס U-הפוך בין גודל המודל לביצועים: רשת מינימלית בת 169 פרמטרים משתווה למודלים גדולים פי 100, בעוד גדלים בינוניים מציגים התאמת-יתר; בתקציב שווה של כ-3,000 פרמטרים הפער בין הארכיטקטורות כמעט נעלם, ממצא המעיד שהקיבולת, ולא הארכיטקטורה, היא הגורם המכריע במשימה זו.

% תרומה משנית מצומצמת בוחנת אפנונים מסדר גבוה: הממצאים הקנוניים עוברים במלואם ל-QPSK מכוח אי-תלות רכיבי I/Q, בעוד 16-QAM חושף רצפת BER מבנית של עיבוד מפוצל לפי ציר; ניסוח הממסר כמסווג דו-ממדי משותף על פני 16 נקודות הקונסטלציה מבטל רצפה זו ומשתווה ל-DF ב-SNR גבוה.

% התרומה המרכזית היא מחקר שיטתי של ממסר מבוסס-למידה תחת ערוצים \emph{בלתי-ידועים ובלתי-מותאמים}. בערוץ ISI בעל שלושה מקדמים, הבלתי-ידוע לשרשרת הקליטה, הממסרים הקלאסיים חסרי-הזיכרון ננעלים על רצפת BER אנליטית של 0.25 : כאשר ביצועי DF אף מחמירים עם עליית ה-SNR : בעוד אותה רשת MLP בת 169 פרמטרים משקמת ממסר אמין עד כדי מרחק של $1$–$1.5$ dB מגלאי Viterbi MLSE בעל ידיעת ערוץ מלאה, וללא כל ידע על משפחת הערוץ. התוצאה נתחמת מכל כיוון: היא לעולם אינה עולה על Viterbi מותאם; על שרשור מורכב של ISI, אי-ליניאריות ופאזה בלתי-ידועה היא נעה כ-$2$ dB מאחורי MLSE מבוסס-פיילוטים ללא כל מודל של הפגם; במשטר נטול-מידע-מוקדם היא משתווה למאזן עיוור (CMA) תוך הימנעות מאי-יציבות של MLSE מבוסס-החלטות; וסריקת תקציב פיילוטים מאתרת סף זיהוי חד הקובע מתי הממסר הלמוד, שאינו נדרש לפיילוטים כלל, עוקף עיבוד קלאסי מבוסס-פיילוטים. לבסוף, ניתוח מורכבות מראה שעלות הממסר הלמוד לכל סימבול קבועה בסדר האפנון ובזיכרון הערוץ (בעוד Viterbi גדל כ-$M^{L}$) ומהירה פי $30$–$90$ בזמן-ריצה. הממסר הלמוד אינו עולה על מקלט קלאסי מותאם, אך תופס נישה מוגדרת היטב של מיתון בלתי-תלוי-משפחה, נטול-זיהוי, וקבוע-מורכבות.

% אסטרטגיית הפריסה המומלצת היא ממסר Hybrid המשלב עיבוד למידה ב-SNR נמוך עם DF ב-SNR גבוה (169 פרמטרים, כ-0.7 KB, פחות מ-3 שניות אימון); עבור 16-QAM, ניסוח כמסווג דו-ממדי משותף; ובכל קישור הנושא פגם בלתי-ממודל (זיכרון, אי-ליניאריות), ממסר למוד מינימלי. התובנה המרכזית היא שממסרים למודים לעולם אינם עולים על האופטימום מודע-המודל, אך ייחודיים בעמידותם כאשר \emph{אין ידיעה איזה מודל חל} : ושמורכבות המודל חייבת להיות מותאמת למבנה המשימה, לא רק לגודלה.

תזה זו מציגה השוואה שיטתית בין אסטרטגיות ממסר קלאסיות ונלמדות לתקשורת דו־דילוגית בערוצי דעיכת ריילי מסוג SISO. נבחנות תשע שיטות: (1) הגבר והעבר (AF) ו־(2) פענח והעבר (DF); (3) רשת פרספטרון רב־שכבתית (MLP) ו־(4) ממסר היברידי מותאם־SNR (Hybrid); (5) מקודד אוטומטי וריאציוני (VAE) ו־(6) רשת יריבותית מותנית (cGAN); ו־(7) Transformer, (8) Mamba-S6 ו־(9) Mamba-2 SSD. הביצועים מוערכים באמצעות סימולציות מונטה קרלו, על בסיס שיעור שגיאות הסיביות (BER), רווחי סמך ומבחני מובהקות סטטיסטית.

בתרחיש ייחוס קנוני שבו מודל הערוץ המשוער תואם לערוץ בפועל, הממסרים הנלמדים הטובים ביותר עולים על AF ביחס אות לרעש (SNR) נמוך ומשיגים ביצועי BER דומים ל־DF ב־SNR בינוני וגבוה, בעוד ש־DF משיגה את שיעור השגיאות הנמוך ביותר מבין השיטות הקונבנציונליות וללא פרמטרים נלמדים. אף שספרות המחקר מבססת תוצאות מידע־תאורטיות עבור DF בתרחישים אסימפטוטיים של קיבולת, תוצאות אלה אינן מעידות במישרין על אופטימליות בתנאים הבלתי־מקודדים ובאורך הסופי הנבחנים כאן. ניתוח סיבוכיות מראה כי הביצועים נקבעים בעיקר על ידי קיבולת המודל ולא על ידי מורכבות הארכיטקטורה: רשת MLP בעלת 169 פרמטרים בלבד משיגה ביצועים המקבילים לרשתות גדולות פי מאה ויותר, ו־Mamba-2 SSD משיגה ביצועים המקבילים ל־Mamba S6 בעלות אימון נמוכה בהרבה.

בהמשך נבחנת עמידות השיטות בתרחישים שבהם הנחות המידול של המקלטים הקונבנציונליים אינן מתקיימות. עבור QPSK מסקנות תרחיש הייחוס נותרות ללא שינוי, ועבור 16-QAM מסווג דו־ממדי משותף מבטל את רצפת השגיאה הנוצרת בעיבוד נפרד של רכיבי המופע והניצב. התרומה המרכזית היא הערכה שיטתית בתנאי ערוץ לא ידועים ובמצבי אי־התאמה בין המודל המשוער לערוץ בפועל, שבהם שיטות קונבנציונליות תלויות־מודל נפגעות במידה ניכרת, בעוד שממסר נלמד בעל 169 פרמטרים בלבד מתקרב לביצועי אלגוריתם ויטרבי לגילוי רצף בנראות מרבית (MLSE) הנהנה מידע מלא על הערוץ, וזאת ללא ידע מפורש על הערוץ.

תוצאות אלה מאפיינות את תחומי הפעולה של הגישות שנבחנו: כאשר המודל המשוער מייצג במדויק את הערוץ, DF משיגה את ביצועי ה־BER הטובים ביותר וללא אימון, ואילו בתנאי אי־ודאות או אי־התאמה הממסרים הנלמדים עמידים יותר לשגיאות מידול. יתרונם העיקרי אינו בעקיפת מקלטים אופטימליים תלויי־מודל בתנאי מידול מדויק, אלא בשימור הביצועים כאשר מודל הערוץ אינו זמין או אינו מתאר כראוי זיכרון ועיוותים.

\end{flushright}
\end{justify}
\end{otherlanguage}
